\documentclass[11pt]{article}
\usepackage[T1]{fontenc}
\usepackage[utf8]{inputenc}
\usepackage[margin=1in]{geometry}
\usepackage{amsmath,amssymb,amsthm}
\usepackage{hyperref}
\hypersetup{hidelinks}
\newtheorem{theorem}{Theorem}[section]
\newtheorem{proposition}[theorem]{Proposition}
\newtheorem{lemma}[theorem]{Lemma}
\newtheorem{corollary}[theorem]{Corollary}
\theoremstyle{remark}
\newtheorem{remark}[theorem]{Remark}
\newcommand{\Q}{\mathbb Q}
\newcommand{\Z}{\mathbb Z}
\newcommand{\F}{\mathbb F}
\newcommand{\C}{\mathbb C}
\newcommand{\R}{\mathbb R}
\newcommand{\PP}{\mathbb P}
\newcommand{\End}{\operatorname{End}}
\newcommand{\Hom}{\operatorname{Hom}}
\newcommand{\Hg}{\operatorname{Hg}}
\newcommand{\Sp}{\operatorname{Sp}}
\newcommand{\SL}{\operatorname{SL}}
\newcommand{\SU}{\operatorname{SU}}
\newcommand{\PSL}{\operatorname{PSL}}
\newcommand{\Res}{\operatorname{Res}}
\newcommand{\Gal}{\operatorname{Gal}}
\newcommand{\NS}{\operatorname{NS}}
\newcommand{\supp}{\operatorname{supp}}
\title{Exact endomorphisms and Hodge classes for a split abelian-cover family\\
\large Revised Phase I: proved scope and general repair boundary}
\author{Kyle Kistner}
\date{4 September 2026}
\begin{document}
\maketitle

\begin{abstract}
For every odd prime \(p\), we give a proof of the generic isogeny
decomposition and endomorphism algebra of the family
\(u^p=x/(x-t),\ v^p=x-1\).
Two local braid eigenvalue ratios determine its independent rank-two
monodromy blocks; explicit involutions supply all multiplicity maps.
Consequently the rational Hodge ring of every power of a very general
Jacobian is generated by divisors.  The proof uses explicitly matched
cyclic-cover results of Menet--Nguyen and classical Hodge and invariant
theory.  It does not use the missing legacy Fox--Gassner certificate or a
Phase II fusion theorem.  We also correct the Galois scope of the moving
part and the orientation of quotient-graph maps.  The former blanket
general finite-abelian block and cross-Hom claims are not asserted here.
This is a mathematical manuscript, not a claim of Lean verification or
external peer review.
\end{abstract}

\section{The restricted theorem and its inputs}
Let \(B=\PP^1_\C\setminus\{0,1,\infty\}\), let \(p\geq3\) be prime, and
let \(X_{p,t}\) denote the smooth projective normalization of
\begin{equation}\label{eq:model}
 u^p=\frac{x}{x-t},\qquad v^p=x-1,\qquad t\in B.
\end{equation}
Write \(J_t=JX_{p,t}\), \(K=\Q(\zeta_p)\), \(F=K^+\), and
\(d=[F:\Q]=(p-1)/2\).  A statement about a very general fibre excludes a
countable union of proper closed subsets of \(B\); it does not describe all
special fibres.

\begin{theorem}\label{thm:main}
For a very general \(t\), there are pairwise non-isogenous absolutely
simple complex abelian varieties \(B_\omega\), indexed by
\(\Omega_p=\F_p^\times/(r\sim r^{-1})\), of dimension \(d\), such that
\begin{align}
 J_t&\sim B_{\{1\}}^2\times B_{\{-1\}}^2
       \times\prod_{\substack{\omega\in\Omega_p\\|\omega|=2}}B_\omega^4,
       \label{eq:decomposition}\\
 \End^0(B_\omega)&=F,\qquad
 \End^0(J_t)=M_2(F)^2\times M_4(F)^{(p-3)/2},\label{eq:end}\\
 \Hg(J_t)^{\rm der}&=
       \bigl(\Res_{F/\Q}\SL_2\bigr)^{(p+1)/2},\label{eq:hg}\\
 \operatorname{rank}\NS(J_t)&=\frac{(p-1)(5p-9)}2.\label{eq:ns}
\end{align}
The genus is \((p-1)^2\).  For every \(q\geq1\), all rational Hodge
classes on \(J_t^q\) are generated by divisor classes and are algebraic.
\end{theorem}

The cyclic-cover input is \cite{MN}, with its standing restriction
\(N\geq3\), \(n\geq3\), finite labels \(1\leq k_i<N\),
\(\gcd(k_1,\ldots,k_n,N)=1\), and primitive
\(q_0=\exp(-2\pi i k/N)\), \(\gcd(k,N)=1\).
Theorem 2.1 gives eigenspace dimension and signature.  Definition 1.1(a)
and Theorem 5.1 give connected special-unitary monodromy when
\(1<\sum_i\{kk_i/N\}<n-1\).
Theorem 2.5 and the reflection calculation immediately before Corollary
2.7 give eigenvalues \(1,q_0^{k_i+k_j}\) in rank two; when the second is
one, Corollary 2.7 gives unipotence.  We do not infer nonidentity from
unipotence.  This input is a preprint theorem, not a theorem proved by
the certificate below.

The other inputs are semisimplicity of geometric polarized variations,
the inclusion of connected monodromy in the very general Hodge group
\cite{Deligne,Andre}, the equivalence between rational weight-one Hodge
maps and isogeny-category maps of complex abelian varieties, the
N\'eron--Severi/Rosati correspondence \cite{BL}, and the first fundamental
theorem for the standard representation of \(\SL_2\) \cite{Weyl}.
The algebraic-group and representation calculations needed here are
given below.

\section{Geometry and algebraic involutions}
Let \(\rho_1,\rho_2\) multiply \(u,v\) by \(\zeta_p\).
The two defining Kummer functions are independent modulo \(p\)-th
powers: their valuations at \(0\) and \(1\) are respectively \((1,0)\)
and \((0,1)\).  Thus the cover is connected of degree \(p^2\), with
group \(\F_p^2\) and inertia of order \(p\) at the four marked points.
Riemann--Hurwitz gives
\[
 2g-2=p^2\bigl(-2+4(1-1/p)\bigr),\qquad g=(p-1)^2.
\]
For \(a,b\in\F_p\), the cyclic quotient is represented by
\(y=u^a v^b\); its four residues are \((a,b,-a,-b)\).
If exactly one of \(a,b\) vanishes, it has two branch points and genus
zero.  If both are nonzero, choose \(1\leq a,b<p\).
The finite labels at \(0,1,t\) can be taken to be
\begin{equation}\label{eq:finite-labels}
 (a,b,p-a).
\end{equation}
Multiplying \(y\) by \(x-t\) accounts for the replacement of \(-a\)
by \(p-a\).  Infinity has residue \(-b\).

Let \(W_{a,b}\) denote the \(q_0=\zeta_p^{-1}\) eigenspace on this
quotient, pulled to \(X\).  As a \(G\)-character it has values
\(\rho_1\mapsto\zeta_p^{-a},\rho_2\mapsto\zeta_p^{-b}\); this fixes the
root-sign convention once and for all.  At every nonzero row \(k\),
\begin{equation}\label{eq:good-match}
 \{ka/p\}+\{kb/p\}+\{k(p-a)/p\}=1+\{kb/p\}\in(1,2).
\end{equation}
Theorem 2.1 of \cite{MN} gives rank two and signature \((1,1)\).
These \(W_{a,b}\), for \(a,b\neq0\), exhaust \(H^1(X,\C)\).
In particular all their Galois conjugates are noncompact factors.

Pass to the finite \'etale base cover \(s^p=t-1\).  On it define
\begin{align}
 \Phi(x,u,v)&=\left(\frac{t(x-1)}{x-t},\frac v s,su\right),\label{eq:phi}\\
 \Psi(x,u,v)&=\left(\frac{t-x}{1-x},-\frac1{su},-\frac s v\right),\label{eq:psi}\\
 \Omega(x,u,v)&=\left(\frac t x,-\frac1v,-\frac1u\right).\label{eq:omega}
\end{align}
Put \(f_1=x/(x-t)\), \(f_2=x-1\), and denote the three first coordinates
by \(T_\Phi,T_\Psi,T_\Omega\).  The identities
\begin{align*}
 f_1(T_\Phi)&=f_2/(t-1),& f_2(T_\Phi)&=(t-1)f_1,\\
 f_1(T_\Psi)f_1&=-1/(t-1),& f_2(T_\Psi)f_2&=1-t,\\
 f_1(T_\Omega)f_2&=-1,& f_2(T_\Omega)f_1&=-1
\end{align*}
prove preservation of the equations, using that \(p\) is odd.
They induce isomorphisms of function fields and hence of the smooth
projective curves.  Direct substitution gives
\[
 \Phi^2=\Psi^2=\Omega^2=1,\qquad \Phi\Psi=\Psi\Phi=\Omega.
\]
Conjugation on the deck group gives, on character labels,
\begin{equation}\label{eq:actions}
 \Phi:(a,b)\mapsto(b,a),\quad
 \Psi:(a,b)\mapsto(-a,-b),\quad
 \Omega:(a,b)\mapsto(-b,-a).
\end{equation}
Their graphs are algebraic correspondences over the base cover.
All monodromy assertions below concern the identity component, which is
unchanged by this finite base change.

\section{Exact rank-two independence from two eigenvalue ratios}
\begin{lemma}\label{lem:projection}
Each \(W_{a,b}\) has connected algebraic monodromy \(\SL_2\) over
\(\C\), and connected projective real Zariski closure \(\PSL_2(\R)\).
\end{lemma}
\begin{proof}
Apply the source input to \(N=p,n=3\) and \eqref{eq:finite-labels}.
Every label is active, the gcd condition holds since \(a\neq0\), and
\eqref{eq:good-match} verifies the good condition with the source's
negative exponential convention.  Its signature is \((1,1)\).
The theorem is initially on the ordered configuration of three finite
points \((b_1,b_2,b_3)\).  Put
\[
 \alpha=b_2-b_1,\qquad \beta=b_1,\qquad
 t=(b_3-b_1)/\alpha,\qquad x=\alpha z+\beta.
\]
This identifies the configuration space with \(\operatorname{Aff}_1\times B\).
On the entire abelian-cover family the equations become
\[
 u^p=\frac{x-b_1}{x-b_3}=\frac{z}{z-t},\qquad
 v^p=x-b_2=\alpha(z-1).
\]
Thus \(u\) is unchanged and \(v=\alpha^{1/p}v_{\mathrm{norm}}\).
Monodromy on the affine-coordinate fibre is a deck transformation,
acting as a finite-order scalar on each character eigenspace.
Consequently the full joint representation into the product of the
projective groups of the \(W_{a,b}\) factors through \(B\).
The same braids \(\alpha_{12},\alpha_{23}\) can therefore be tested
before or after normalization, simultaneously for every pair of factors.
In particular the connected projective images are preserved.
The determinants of braid generators are \(p\)-th roots of unity by
Corollary 2.7; the same holds after normalization.  Thus the connected
linear image is \(\SU(1,1)\), whose complexification is \(\SL_2\).
\end{proof}

\begin{lemma}[Normalizer of a graph]\label{lem:normalizer}
Let \(H\) be a connected centreless simple algebraic group in
characteristic zero.  For an automorphism \(\theta\) of \(H\), the
normalizer in \(H\times H\) of
\(\Gamma_\theta=\{(h,\theta(h)):h\in H\}\) is \(\Gamma_\theta\).
If two projective monodromy factors form a connected graph, their full
projective representations obey that graph relation.
\end{lemma}
\begin{proof}
If \((a,b)\) normalizes the graph, then for every \(h\),
\(b\theta(h)b^{-1}=\theta(aha^{-1})\).  Therefore
\(\theta(a)^{-1}b\) centralizes \(H\) and is the identity.
For the second assertion, the identity component of an algebraic
closure is normal in the full closure; apply the first assertion to
the two-coordinate closure.  In particular one must not raise a
finite-order braid image to a power and then infer a graph obstruction
from the resulting identity.
\end{proof}

\begin{lemma}\label{lem:two-ratios}
Two split character factors can lie in one connected projective
monodromy graph only if
\begin{equation}\label{eq:orbit}
 (a',b')\in\{(a,b),(-a,-b),(b,a),(-b,-a)\}.
\end{equation}
\end{lemma}
\begin{proof}
Use the same ordered finite labels \(0,1,t\) for both factors.
For the pure braids \(\alpha_{12}\) and \(\alpha_{23}\), the source
reflection formula gives unordered eigenvalue ratios
\[
 \{\zeta_p^{a+b},\zeta_p^{-(a+b)}\},\qquad
 \{\zeta_p^{b-a},\zeta_p^{-(b-a)}\}.
\]
The displayed signs absorb both the fixed negative-root convention
and the ordering of the two eigenvalues.  This assertion also holds
when a ratio is one; no nonidentity assertion is used.
By Lemma \ref{lem:normalizer} a connected graph forces equality of
these projective conjugacy invariants.  Algebraic automorphisms of
\(\PSL_2\) over \(\C\) are inner, so they preserve these invariants.
For independent signs \(\epsilon,\eta\in\{1,-1\}\), we obtain
\[
 a'+b'=\epsilon(a+b),\qquad a'-b'=\eta(a-b)\quad\text{in }\F_p.
\]
As \(2\) is invertible,
\[
 a'=\frac{\epsilon+\eta}{2}a+\frac{\epsilon-\eta}{2}b,\qquad
 b'=\frac{\epsilon-\eta}{2}a+\frac{\epsilon+\eta}{2}b.
\]
Equal signs give \((a',b')=\epsilon(a,b)\), and opposite signs give
\((a',b')=\epsilon(b,a)\).  This includes \(a=b\) and \(a=-b\).
\end{proof}

\begin{proposition}\label{prop:blocks}
The connected complex monodromy is a product of standard \(\SL_2\)
factors, one for each orbit \eqref{eq:orbit}.  Each such factor acts
diagonally on the rank-two summands in that orbit and trivially on the
other orbits.  There are \((p^2-1)/4\) factors: \(p-1\) orbits of size
two and \((p-1)(p-3)/4\) orbits of size four.
\end{proposition}
\begin{proof}
Semisimplicity and the surjections in Lemma \ref{lem:projection} imply
the usual semisimple Goursat decomposition into graph blocks.
Indeed each simple Lie ideal projects either to zero or isomorphically
to a simple target; two different ideals cannot both project
nontrivially to one target, since their images commute.
Lemma \ref{lem:two-ratios} excludes graphs between distinct displayed
orbits.  Conversely \eqref{eq:actions} supplies isomorphisms of local
systems throughout each orbit after the finite base change, forcing
the same monodromy action there.  Connected determinants are trivial;
lifting a projective graph gives the diagonal \(\SL_2\) on the
standard modules.  The standard two-dimensional representation is
faithful, so no additional quotient or connected scalar factor remains.

Only the lines \(a=b\) and \(a=-b\) have size-two orbits, giving \(p-1\)
such orbits.  All other orbits have size four.  Equivalently Burnside
gives \(((p-1)^2+2(p-1))/4=(p^2-1)/4\).
\end{proof}

\begin{remark}
This proof uses neither an exact Fricke trace product nor a
half-shift exclusion at even denominators, and does not depend on
the missing Fox--Gassner source.  It proves this split odd-prime case;
it is not a replacement for those calculations for arbitrary disk words.
\end{remark}

\section{Algebraic matrix units and the rational commutant}
\begin{lemma}[Use inverse maps, not an oriented Rosati shortcut]\label{lem:units}
Suppose algebraic correspondences give isomorphisms
\(T_i:W_0\longrightarrow W_i\), with algebraic inverses
\(S_i:W_i\longrightarrow W_0\).  Include the character projectors in
these maps, so all maps vanish off their stated source.  Then
\(E_{ij}=T_iS_j\) satisfy \(E_{ij}E_{k\ell}=\delta_{jk}E_{i\ell}\).
\end{lemma}
\begin{proof}
The middle composite \(S_jT_k\) is zero for \(j\neq k\), and is
\(1_{W_0}\) for \(j=k\).  All compositions remain algebraic.
For an automorphism of \(X\), its inverse graph gives \(S_i\).
For a quotient \(q:X\to Y\), the identity
\(q_*q^*=(\deg q)1\) gives the required rational normalization.
No square roots of Schur scalars are needed.
\end{proof}

\begin{remark}
For the bilinear polarization, the cyclotomic-linear extension of
Rosati satisfies \(e_c^\dagger=e_{-c}\).  The adjoint of an isolated
arrow \(W_c\to W_e\) has source \(W_{-e}\), not in general \(W_e\).
Thus the old formula using \(\Theta_i^\dagger\Theta_i\) on an oriented
corner cannot justify these matrix units.  Explicit inverse graphs do.
\end{remark}

Rational deck-isotypic factors are indexed by
\(L_r=\F_p(1,r)\), \(r\in\F_p^\times\).  Denote their abelian
varieties by \(A_r\); their dimension is \(p-1\), since each has
\(p-1\) complex eigenspaces of rank two.  The rational deck algebra
on \(A_r\) is \(K\).  The algebraic involution \(\Psi\) preserves
\(A_r\), has square one, and conjugates \(K\) by complex conjugation.
Hence there is a unital algebra map
\[
 K\rtimes\Gal(K/F)\longrightarrow\End^0(A_r).
\]
The crossed product is \(\End_F(K)=M_2(F)\): multiplication by
\(K\) and conjugation act faithfully on the two-dimensional
\(F\)-space \(K\), and both algebras have dimension four over \(F\).
Its simplicity makes the unital map injective.  A primitive idempotent
gives \(A_r\sim B_r^2\), with \(\dim B_r=d\).
The map \(\Phi\) identifies \(A_r\) with \(A_{r^{-1}}\).
For each nonfixed pair, its graph and inverse combine the two
\(M_2(F)\) actions into \(M_4(F)\); possible twists of the identified
copy of \(F\) do not change this algebra isomorphism.
Thus there is an algebraic subalgebra
\begin{equation}\label{eq:lower-algebra}
 \mathcal A=M_2(F)^2\times M_4(F)^{(p-3)/2}
       \ \subseteq\ \End^0(J_t).
\end{equation}
All rational idempotents and inverses here are in the isogeny category.

\begin{proposition}\label{prop:commutant}
For the connected monodromy group \(M\),
\(\End_M(H^1(J_t,\Q))=\mathcal A\) (with the harmless opposite
algebra convention for pullback on \(H^1\)).
\end{proposition}
\begin{proof}
Over \(\C\), Proposition \ref{prop:blocks} and Schur's lemma give
exactly one matrix algebra \(M_m(\C)\) for each oriented orbit of
size \(m=2\) or \(4\), and no maps between different orbits.
Lemma \ref{lem:units}, applied to \(\Phi,\Psi,\Omega\), supplies all
of these matrix units algebraically.  Equivalently the complex
dimension of the commutant is
\[
 4(p-1)+16\frac{(p-1)(p-3)}4=4(p-1)(p-2).
\]
This equals the rational dimension of \(\mathcal A\) in
\eqref{eq:lower-algebra}.  The algebraic correspondences commute
with connected monodromy after the finite base change.  Equality
after scalar extension, or equality of these dimensions and the
inclusion, proves the rational assertion.  No unsupported
descent of a chosen single cyclotomic arrow is used.
\end{proof}

For a very general fibre the connected monodromy group lies in the
Hodge group, so
\[
 \mathcal A\subseteq\End_{\rm HS}(H^1)
       \subseteq\End_M(H^1)=\mathcal A.
\]
Weight-one full faithfulness gives \eqref{eq:end}.
Its primitive idempotents give \eqref{eq:decomposition}.
Different \(\omega\) have no Hom corner and are not isogenous;
\(\End^0(B_\omega)=F\) has no nontrivial idempotent, so \(B_\omega\)
is simple, and hence absolutely simple over \(\C\).

The centralizer of \(\mathcal A\) in the rational symplectic group is
\[
 C_{\Sp}(\mathcal A)=
       \prod_{\omega\in\Omega_p}\Res_{F/\Q}\Sp_2
       =\prod_{\omega\in\Omega_p}\Res_{F/\Q}\SL_2.
\]
To check its form, on one component write
\(H^1=F^2\otimes_F F^{m_\omega}\).  Positivity of Rosati fixes the
totally real centre \(F\) and makes its involution on the split matrix
algebra \(M_{m_\omega}(F)\) orthogonal: it is adjoint to a nondegenerate
symmetric multiplicity form.  The polarization is therefore the trace
of the product of that form with an alternating form on \(F^2\).
The operators commuting with \(M_{m_\omega}(F)\) therefore preserve
this nondegenerate rank-two alternating form.  Every such form has a
symplectic basis over \(F\), so its group is the split \(\Sp_2=\SL_2\).
This centralizer is connected.
By Proposition \ref{prop:blocks}, \(M\) has the same complexification
and dimension as this centralizer, hence equals it over \(\Q\).
Algebraicity of \(\mathcal A\) now gives
\[
 M\subseteq\Hg\subseteq C_{\Sp}(\mathcal A)=M.
\]
In particular \eqref{eq:hg} holds; this argument also shows that the
Hodge group itself is semisimple in this restricted family.

On \(M_m(F)\), a positive Rosati involution is adjoint with respect
to a symmetric form over \(F\).  Its fixed subspace has dimension
\(d\,m(m+1)/2\) over \(\Q\).
The N\'eron--Severi/Rosati correspondence therefore gives
\[
 \operatorname{rank}\NS(J_t)
 =d\left(2\cdot3+\frac{p-3}{2}\cdot10\right)
 =\frac{(p-1)(5p-9)}2.
\]

\section{All powers}
\begin{proof}[Proof of the last assertion of Theorem \ref{thm:main}]
Fix \(q\geq1\).  After scalar extension, \(H^1(J_t^q)\) is a direct
sum of standard two-dimensional modules for the independent
\(\SL_2\) factors, with their multiplicities multiplied by \(q\).
The first fundamental theorem for \(\SL_2\) states that tensor
invariants are spanned by products of alternating pair contractions.
Passing to the exterior algebra by antisymmetrization shows that its
invariant algebra is generated by invariant classes of degree two.

Every multiplicity map is algebraic by Proposition
\ref{prop:commutant}, including the matrix maps between the \(q\)
copies.  Together with a polarization these give the invariant
bilinear pairings.  Their alternating parts are algebraic divisor
classes, equivalently the corresponding Poincar\'e classes on the
products and their pullbacks.  Thus the entire \(M\)-invariant
exterior algebra is generated by divisors.
Since \(M\subseteq\Hg\), all rational Hodge classes are among these
invariants.  Conversely divisors are Hodge classes.  Equality follows
over \(\C\) and then over \(\Q\) by faithful scalar extension.
This uses no central-torus or fusion theorem.
\end{proof}

\section{Correct moving parts for general abelian branch data}
This section repairs a general reduction, without asserting a new
complete general block theorem.
Let a full family of connected finite abelian covers of \(\PP^1\)
have rational cohomology \(V\) and connected monodromy \(M\).
Its deck algebra decomposes \(V=\bigoplus_{\mathcal O}V_{\mathcal O}\)
into rational character orbits.  An orbit means a full cyclotomic
Galois orbit, not just the pair \(\{c,-c\}\).
For a character with \(s\) nonzero branch residues
\(0<\alpha_i<1\), set \(Q(c)=\sum_i\alpha_i\in\Z\).
Its two Hodge multiplicities, in one of the two orientations, are
\(s-1-Q(c)\) and \(Q(c)-1\).
Call the orbit positive if some character in it has both
multiplicities positive.

\begin{proposition}[Galois-saturated moving criterion]\label{prop:moving}
Assume that every positive character in the family has full connected
factorwise monodromy on its eigenspace, and hence no invariant vector.
Then, after a finite \'etale base change,
\[
 V^M=\bigoplus_{\mathcal O\ {\rm not\ positive}}V_{\mathcal O},
 \qquad
 V_{\rm mov}=(V^M)^\perp
      =\bigoplus_{\mathcal O\ {\rm positive}}V_{\mathcal O}.
\]
All complex embeddings of a positive rational orbit must be retained,
including its definite compact embeddings.
\end{proposition}
\begin{proof}
The \(M\)-fixed subspace is rational and stable under the deck field
of each \(V_{\mathcal O}\).  At a positive embedding its scalar
extension is zero by the stated projection premise.  A vector space
over that deck field has equal dimension at all embeddings, so this
fixed subspace is zero on the whole rational orbit.

If no character in an orbit is positive, every real character pair
has definite Hodge signature.  The real deck centralizer preserving
the polarization is then a product of compact unitary groups.
The whole rational summand, unlike an isolated complex eigenspace,
has a monodromy-stable integral lattice, obtained by intersection
with the integral cohomology lattice.  Its monodromy lies in a
compact group and in a discrete integral matrix group, hence is
finite.  A common finite \'etale cover kills these finitely many
finite images.  The deck decomposition is orthogonal between
rational orbits, proving the stated orthogonal-complement formula.
\end{proof}

The projection premise has a precise source match for characters of
effective order \(N\geq3\) on the full branch family.  Put one active
branch point at infinity.  The \(n=s-1\) finite weights have sum
\(Q(c)-\alpha_\infty\).  Positivity gives
\[
 1<Q(c)-\alpha_\infty<s-2=n-1.
\]
Their integral labels are nonzero, the effective-order condition gives
the gcd equal to one, and either primitive orientation satisfies the
same strict inequalities.  These are the hypotheses of
\cite[Definition 1.1(a), Theorem 5.1]{MN}.
Forgetting inactive labels is a configuration-space fibration with
connected fibre and surjects on fundamental groups.  Normalization
of the active labels changes only finite scalar monodromy, as above.
Effective order two uses the separate hyperelliptic symplectic
monodromy input of A'Campo \cite{ACampo}, not Theorem 5.1.
Thus the proposition is an exact general representation-theoretic
repair; it does not by itself determine graph identifications between
different rational orbits.

The old lemma discarding each definite complex eigenspace is false.
For example \(c=(1,1,1,2)/5\) has signatures, over its four embeddings,
\((2,0),(1,1),(1,1),(0,2)\).  The second embedding has full noncompact
monodromy.  If the first had finite image, applying the cyclotomic
field automorphism to a basis over \(\Q(\zeta_5)\) would make the
second image finite too.  This contradicts full projection.

\section{Correctly oriented general quotient arrows and boundaries}
For four active labels, let \(P\) be a double transposition, and let
cyclic words \(c,e\) have the same effective order \(N\).
Fix the same primitive root \(q_0\) when defining the eigenspaces
on \(Y_c:y^N=f_c(x)\) and \(Y_e:v^N=f_e(x)\).
If \(e=Pc\), the divisor identity
\[
 f_e(T_P(x))/f_c(x)=a(t)h(x)^N
\]
gives, after adjoining a root of the unit \(a(t)\), an algebraic
deck-equivariant map
\[
 (x,y)\longmapsto\bigl(T_P(x),a(t)^{1/N}h(x)y\bigr).
\]
It identifies \(W_c\) with \(W_e\).  The unit is a constant times
integral powers of \(t,t-1\), from the displayed double-transposition
M\"obius functions, so the base change is \'etale over \(B\).

If \(e=-Pc\), the analogous map has second coordinate proportional
to \(h(x)/y\).  It conjugates the cyclic generator to its inverse,
so the \(q_0\) eigenspace on \(Y_c\) maps to the \(q_0^{-1}\)
eigenspace on \(Y_e\), that is to \(W_{-e}\), not automatically
\(W_e\).  Hence geometric oriented arrows are indexed by the
actual relations \(e=Pc\).  A relation only up to sign classifies
possible projective equivalence, not a specified Hodge Hom corner.

For a genuine oriented arrow, pull--graph--push and the normalization
\((\deg q_c)^{-1}q_c{}_*\) give an isomorphism on \(W_c\).
The inverse graph and \((\deg q_e)^{-1}q_e{}_*\) give its inverse.
Lemma \ref{lem:units} therefore proves matrix-algebra saturation on
each actually connected oriented graph component.  Taking the
Galois-stable algebra of these correspondences permits rational
descent, but supplies no arrows between disconnected oriented
components.

In particular, a missing internal standard--dual map can be transported
across an existing graph to a missing cross-coordinate map.  It is
incorrect to exclude internal corners from a theorem and then assert
that every off-diagonal \(M\)-intertwiner is a graph Hodge map.
In \eqref{eq:model}, the explicit \(\Psi\) supplies precisely the
needed orientation maps, which is why the restricted theorem is
unaffected.

The following former claims are withdrawn as unconditional statements
in this revision: the real monodromy product obtained by deleting all
definite complex eigenspaces; the blanket equality of graph algebras
with all cross-character \(\End_M\) corners; and the complete
finite-abelian Hodge-group or all-powers theorem.  Reinstating a general
result requires an exact block theorem on the corrected rational
representation (including compact embeddings), correctly oriented
algebraic generators, and a proof of the remaining determinant-tensor
algebraicity.  General support separation, high-rank reconstruction,
and the even-denominator four-point fingerprint argument retain their
own proof/source obligations; none is silently treated as an external
theorem here.  Phase II's global fusion and cycle-specialization
obligations remain separate.

\appendix
\section{Reproducible checks and provenance}
The new \path{scripts/check_phase_i.py} is a standalone symbolic and
finite regression script using SymPy.  It checks the six rational
identities, all involution and deck-conjugation relations, the
two-eigenratio fibres, the source inequalities, and the orbit and
dimension formulas for every odd prime up to a configurable bound
(default \(101\)).  It includes small exact crossed-product and
inverse-matrix-unit checks and records the new source hashes.
These finite checks are not the proof for all primes, nor a
formalization of monodromy or Hodge theory.

The original manuscripts and \path{legacy/certificates/} are unchanged
provenance copies.  In particular the absent source
\path{phase1src/p_ary_mt_defect_certificate.py} has not been
reconstructed or represented as recovered.  The frozen
\(7{,}077{,}120\)-vector and Fox--Gassner claims are not outputs of the
new script.  The restricted proof above no longer needs them.

\begingroup
\small
\begin{thebibliography}{9}
\setlength{\itemsep}{3pt}
\bibitem{MN} G. Menet and D.-M. Nguyen,
\emph{Representations of braid groups via cyclic covers of the sphere:
Zariski closure and arithmeticity}, arXiv:2310.10401v3 (2024),
\url{https://arxiv.org/html/2310.10401v3}.
\bibitem{ACampo} N. A'Campo,
\emph{Tresses, monodromie et le groupe symplectique},
Comment. Math. Helv. 54 (1979), 318--327.
\bibitem{Deligne} P. Deligne, \emph{Th\'eorie de Hodge II},
Publ. Math. IH\'ES 40 (1971), 5--57.
\bibitem{Andre} Y. Andr\'e,
\emph{Mumford--Tate groups of mixed Hodge structures and the theorem of
the fixed part}, Compositio Math. 82 (1992), 1--24.
\bibitem{BL} C. Birkenhake and H. Lange,
\emph{Complex Abelian Varieties}, second edition, Springer, 2004.
\bibitem{Weyl} H. Weyl,
\emph{The Classical Groups: Their Invariants and Representations},
Princeton University Press, 1939.
\end{thebibliography}
\endgroup
\end{document}
